% CODASPY 2027 draft. Build: latexmk -pdf main.tex   (after: python make_numbers.py)
% Submission uses [sigconf,review,anonymous]; camera-ready drops review,anonymous.
\documentclass[sigconf,review,anonymous]{acmart}

\usepackage{xspace}
\usepackage{booktabs}
\usepackage{listings}
% Generated by paper/make_numbers.py from evaluations/results and conference-snapshot. Do not edit.
\newcommand{\NumTOneFixtures}{21\xspace}
\newcommand{\NumTOnePathTP}{14\xspace}
\newcommand{\NumTOnePathFP}{0\xspace}
\newcommand{\NumTOnePathFN}{0\xspace}
\newcommand{\NumTOnePathPrecision}{100.0\%\xspace}
\newcommand{\NumTOnePathRecall}{100.0\%\xspace}
\newcommand{\NumTOneMedianMs}{597\xspace}
\newcommand{\NumTOnePNinetyFiveMs}{820\xspace}
\newcommand{\NumTOneCheckov}{354\xspace}
\newcommand{\NumTOneTfsec}{282\xspace}
\newcommand{\NumTOneTrivy}{222\xspace}
\newcommand{\NumTOneHopCheckov}{28.6\%\xspace}
\newcommand{\NumTOneHopTfsec}{59.2\%\xspace}
\newcommand{\NumTOneHopTrivy}{34.7\%\xspace}
\newcommand{\NumTOnePytmOnly}{1682\xspace}
\newcommand{\NumTOnePytmFed}{1609\xspace}
\newcommand{\NumTOnePytmDecided}{94\xspace}
\newcommand{\NumTOnePytmAppTeam}{1515\xspace}
\newcommand{\NumTOnePytmRuledOut}{176\xspace}
\newcommand{\NumTOnePackFindings}{61\xspace}
\newcommand{\NumPackThreats}{7\xspace}
\newcommand{\NumTOneAIGenerated}{70\xspace}
\newcommand{\NumTOneAIAccepted}{65\xspace}
\newcommand{\NumTOneAIHeld}{5\xspace}
\newcommand{\NumTOneRemPaths}{14\xspace}
\newcommand{\NumTOneRemBreak}{10\xspace}
\newcommand{\NumTOneRemRejected}{2\xspace}
\newcommand{\NumTOneRemWhatIfAgree}{17\xspace}
\newcommand{\NumTOneRemWhatIfChecked}{17\xspace}
\newcommand{\NumTTwoCases}{10\xspace}
\newcommand{\NumTTwoTP}{33\xspace}
\newcommand{\NumTTwoFN}{3\xspace}
\newcommand{\NumTTwoFP}{1\xspace}
\newcommand{\NumTTwoTN}{4\xspace}
\newcommand{\NumTTwoPrecision}{97.1\%\xspace}
\newcommand{\NumTTwoRecall}{91.7\%\xspace}
\newcommand{\NumTBVulns}{13\xspace}
\newcommand{\NumTBSamples}{9\xspace}
\newcommand{\NumTBScoredSamples}{3\xspace}
\newcommand{\NumTBTerraraThreatsAll}{715\xspace}
\newcommand{\NumTBTerraraThreats}{488\xspace}
\newcommand{\NumTBTerraraCovered}{13\xspace}
\newcommand{\NumTBTerraraPerVuln}{37.5\xspace}
\newcommand{\NumTBPath}{5\xspace}
\newcommand{\NumTBAsked}{8\xspace}
\newcommand{\NumTBConfirmed}{0\xspace}
\newcommand{\NumTBNone}{0\xspace}
\newcommand{\NumTBStrict}{5\xspace}
\newcommand{\NumTBDeveloper}{13\xspace}
\newcommand{\NumTBFirstScreen}{8\xspace}
\newcommand{\NumTBItems}{171\xspace}
\newcommand{\NumTBItemsFirst}{31\xspace}
\newcommand{\NumTBItemsPerVuln}{13.2\xspace}
\newcommand{\NumTBScannerFindings}{617\xspace}
\newcommand{\NumTBFirstRunCovered}{6\xspace}
\newcommand{\NumTBFirstRunPaths}{0\xspace}
\newcommand{\NumDemoOpen}{110\xspace}
\newcommand{\NumDemoSixSettle}{34\xspace}
\newcommand{\NumDemoQuestions}{37\xspace}
\newcommand{\NumDemoGeneric}{122\xspace}
\newcommand{\NumDemoPublished}{118\xspace}
\newcommand{\NumDemoConfirmed}{8\xspace}
\newcommand{\NumDemoRuledOut}{12\xspace}
\newcommand{\NumDemoAIGenerated}{5\xspace}
\newcommand{\NumDemoAIAccepted}{5\xspace}


\newcommand{\sys}{\textsc{PlanPath}\xspace}   % anonymized tool name for review; see README.md
\newcommand{\todo}[1]{\textcolor{red}{[TODO: #1]}}
\newcommand{\conf}{\textsf{confirmed}\xspace}
\newcommand{\ask}{\textsf{app-team}\xspace}
\newcommand{\ruled}{\textsf{ruled-out}\xspace}
\setlength{\emergencystretch}{2em}
\lstset{basicstyle=\ttfamily\footnotesize,breaklines=true,columns=fullflexible,frame=single}

\setcopyright{none}
\acmConference[CODASPY '27]{Seventeenth ACM Conference on Data and Application Security and Privacy}{June 14--17, 2027}{Fort Collins, CO, USA}
\acmBooktitle{CODASPY '27}

\begin{document}

\title{From Terraform Plans to Attack Paths:\\Evidence-Fed Threat Modeling that Developers Can Answer}

\author{Anonymous Author(s)}
\affiliation{\institution{Anonymous Institution}\country{}}

\begin{abstract}
Infrastructure as code (IaC) put the cloud architecture of most applications into files that developers review every day,
yet threat modeling still happens elsewhere, if at all. IaC scanners report isolated misconfigurations; threat modeling
tools that read IaC report every threat an architecture could have. Neither tells a developer which few changes in
\emph{this} pull request let an outsider reach sensitive data, nor asks them the handful of questions only they can answer.

We present \sys, which turns a Terraform plan into (i)~\emph{attack paths}: chains of network and IAM relationships from an
attacker-controlled entry to a classified asset, where every hop cites plan or source evidence and every assumption is named;
(ii)~an \emph{evidence-fed} threat model: an unmodified OWASP pytm threat library whose element attributes are set from the
plan, with each threat classified as \emph{confirmed} by the plan, \emph{ruled out} by it, or \emph{open}; and (iii)~a few
plain yes/no questions, selected greedily, that settle the open threats on the paths first. Answers become attested evidence
that can never change a path. Language models only re-tell what the evidence establishes, under a validator, or propose a patch
that the pipeline re-plans and proves.

On \NumTOneFixtures synthetic fixtures \sys recovers every labeled path (\NumTOnePathTP/\NumTOnePathTP); on \NumTTwoCases
public vulnerable-by-design cases its reachability precision and recall are \NumTTwoPrecision and \NumTTwoRecall. On the
samples of TerrARA's relevance benchmark that document their vulnerabilities, it points at all \NumTBDeveloper of
\NumTBVulns known vulnerabilities with \NumTBItems items to read, where TerrARA reports \NumTBTerraraThreats threats;
\NumTBPath of them as attack paths that match the published walkthrough. \todo{developer study result, one sentence.}
\end{abstract}

\maketitle

%======================================================================
\section{Introduction}
%======================================================================

A pull request that changes Terraform is an architecture change. It can open a security group, attach a role to a server,
widen an IAM policy with a wildcard, or put a public API in front of a function. Each of these is a line a reviewer can see;
whether they \emph{combine} into a way in is not. That combination is what threat modeling is for, and it is the part
developers are least equipped to do: few have been trained in it, and the artifacts it produces---data flow diagrams,
STRIDE tables, lists of threats---are written for security specialists~\cite{shostack2014threat,tahaei2019survey}.

Two families of tools read IaC today. \emph{Misconfiguration scanners} (Checkov, tfsec, Trivy) check resources against
rules and report findings one resource at a time; on our synthetic corpus they emit \NumTOneCheckov, \NumTOneTfsec and
\NumTOneTrivy findings, and at best \NumTOneHopTfsec of the hops on the labeled attack paths carry a finding at all.
\emph{Threat modeling tools} extract an architecture from IaC and apply a threat library to it. TerrARA~\cite{tran2025terrara},
the closest published work, builds an enriched DFD from Terraform and elicits threats with SPARTA; on its relevance benchmark it
covers 26 of 28 known vulnerabilities with 715 threats. Architecture-only threat models are necessarily generic: without the
plan's values they cannot tell a wildcard grant from a scoped one, or an instance that requires IMDSv2 from one that does not,
so they must report every threat the architecture \emph{could} have. The same is true of pytm~\cite{pytm} run on a model
generated from the plan: \NumTOnePytmOnly threats on our corpus, none of them connected into a path.

We argue that the plan already contains much of what these tools leave open, and that what it does not contain can be
reduced to a few questions a developer can answer. Concretely:

\begin{itemize}
\item \textbf{Attack paths with evidence (\S\ref{sec:paths}).} A typed security graph built from the plan's values
  \emph{and} its configuration references, with network and IAM semantics evaluated deterministically. An attack path is a
  chain of relationships from an attacker-controlled entry to a classified asset; every hop cites evidence records that point
  at a file and line, and every hop that needs more than the configuration shows carries a named precondition.
\item \textbf{Evidence-fed threat modeling (\S\ref{sec:classify}).} We generate an executable pytm model from the same graph,
  set element attributes from plan evidence, and classify each (threat, element) pair of pytm's unmodified library as
  \conf by the plan, \ruled by it, or \ask---open until someone who knows the application answers.
\item \textbf{Questions instead of threat lists (\S\ref{sec:questions}).} The open threats are reduced to plain yes/no
  questions on (element, control) pairs, chosen greedily and path-first. Answers are committed next to the Terraform, become
  attested evidence, re-decide the threats, and by construction never add, remove or change an attack path.
\item \textbf{Output for people who do not threat-model (\S\ref{sec:people}).} Each path is also rendered as a one-screen
  card in plain language, for a developer (with the line to change) or for anyone. Language models may re-tell a path or
  propose a patch, but an annotation is accepted only if a validator finds every claim backed by cited evidence, and a patch
  only if the pipeline re-plans it and proves its effect.
\item \textbf{Evaluation (\S\ref{sec:eval}).} Path accuracy on \NumTOneFixtures synthetic fixtures and \NumTTwoCases public
  cases; a head-to-head on TerrARA's benchmark with a ground truth committed before the first run; reviewer load; and the
  protocol of a developer comprehension study. \todo{study results.}
\end{itemize}

%======================================================================
\section{Motivating example}\label{sec:example}
%======================================================================

Figure~\ref{fig:card} shows what \sys reports for a 23-resource plan: an internet-facing load balancer forwards to an EC2
instance whose metadata service accepts IMDSv1; the instance profile's role has \texttt{s3:GetObject} on
\texttt{arn:aws:s3:::*}; one bucket is tagged confidential. Every scanner flags some of these resources. None says that they
form one path, that the wildcard is what makes the path reach \emph{that} bucket, or that requiring IMDSv2 only lowers its
priority while scoping the grant removes it. pytm on the architecture alone reports \NumDemoGeneric threats for this
plan. Fed with the plan, \NumDemoConfirmed are confirmed (one on every hop of the path), \NumDemoRuledOut are ruled out,
and six yes/no questions settle \NumDemoSixSettle of the \NumDemoOpen that remain open.

\begin{figure*}[t]
\begin{lstlisting}
HIGH: Someone on the internet could list and read the customer-data bucket
(aws_s3_bucket.customer_data), which holds confidential data, if they find a
way to make the app server (aws_instance.app) fetch a web address for them
(an SSRF bug) or run code on it.
 1. Anyone on the internet can reach the web load balancer (HTTPS/443).  alb.tf:3
 2. The web load balancer passes those requests on to the app server.  network.tf:94-100
 3. The app server can fetch the cloud keys of the app role from a local
    address (IMDS). Because the older IMDS setting is allowed, a request
    someone tricks the server into making can fetch them too.
    (only if the condition above holds)  app.tf:10-13
 4. Those keys are allowed to list and read the customer-data bucket because
    the permission covers everything of that kind (a wildcard).  iam.tf:25-35
 Fix: Narrow the permission of the app role to the resources it needs.
    The path engine re-ran itself with this change: the path disappears.
 ? Does the customer-data bucket hold personal data?
 ? Does the app server check every request's input before using it?
\end{lstlisting}
\caption{The developer card for the demo plan (abridged). Every sentence is generated from evidence records.}
\label{fig:card}
\end{figure*}

%======================================================================
\section{From plan to attack paths}\label{sec:paths}
%======================================================================

\paragraph{Input.} \texttt{terraform show -json} on a saved plan gives planned values, unknown-after-apply markers, and the
configuration's unevaluated expressions with their references~\cite{tf-json}. Plans are produced offline with fake
credentials, a provider filesystem mirror, and an AWS API mock for data sources; no cloud account is involved.

\paragraph{Security graph.} Each resource becomes a node (modeled or not, so coverage stays measurable). Edges are typed:
\texttt{EXPOSES} (internet to load balancer, public instance, public database, unauthenticated function URL or API route),
\texttt{ROUTES\_TO} (listener to target group to instance or ECS service; API integration to backend), \texttt{ASSUMES}
(workload to role via instance profile, execution role, task role, or container host; principal to role via trust or a
privilege-escalation primitive), \texttt{ALLOWS} (principal to asset with an access class), and \texttt{ESCALATES}
(principal to account administrator). IAM evaluation implements the stated subset of AWS policy
logic~\cite{aws-iam-eval}: explicit deny, NotAction/NotResource, permission boundaries, resource policies, conditions (kept
as preconditions, not evaluated), and a catalogue of escalation primitives~\cite{rhino-privesc,iamvulnerable}. Policies
whose documents are unknown at plan time are recovered \emph{symbolically} from source, replacing interpolations such as
\texttt{\$\{aws\_s3\_bucket.x.arn\}} by tokens that match only that resource.

\paragraph{Attack path.} \emph{An attack path is a sequence of typed graph relationships connecting an attacker-controlled
entry condition to a security-relevant asset, where every edge is supported by deterministic evidence and any
non-deterministic assumption is explicitly labeled.} Entries are the internet and, optionally, IAM users (credential
compromise). Assets are data stores classified by tag or type, and the account administrator. A hop is a \emph{fact} when the
edge suffices, an \emph{assumption} when crossing it needs a named precondition: SSRF or code execution in the
workload, an escape from a container, stolen user credentials, a satisfied IAM condition, or traffic from an allowed public
range. Priority is a documented integer score of target class, entry and preconditions.

\paragraph{What-if.} For each hop, the corresponding fix is applied to a copy of the graph and the search re-run; the report
says whether the change breaks the path, downgrades it, or leaves an alternative route.

\paragraph{Evidence ledger.} Every edge, fact and path refers to ledger records (\texttt{TF-}, \texttt{NET-}, \texttt{IAM-},
\texttt{GRAPH-}, \texttt{PATH-}, \texttt{SCAN-}, \texttt{APP-}) with a source location. Record IDs are assigned after the run in
sorted order, so the same plan yields byte-identical artifacts.

%======================================================================
\section{Evidence-fed threat classification}\label{sec:classify}
%======================================================================

pytm represents a threat $t$ as a target class and a condition $\varphi_t$, a Python expression over attributes of the element
(or of a data flow's source, sink and data). Attributes default to \texttt{False}; a finding is reported when $\varphi_t$
evaluates true. \sys generates a pytm model from the graph (servers, functions, datastores, flows, data objects, boundaries,
and the IMDS credential hop) and sets attributes from evidence, e.g.\ \texttt{isEncrypted} from a listener's protocol,
\texttt{hasAccessControl} from a bucket policy, \texttt{authenticatesSource} from IMDS settings.

\paragraph{Formalization.} Fix an element $e$ and threat $t$. Let $V_t$ be the attributes $\varphi_t$ reads. The plan
decides $D \subseteq V_t$ with values $\sigma: D \to \mathit{Val}$, each value backed by evidence; $U = V_t \setminus D$ is
undecided. Write $\varphi_t(\sigma, \tau)$ for the condition under $\sigma$ and an assignment $\tau$ to $U$, and $\mathbf{0}$
for pytm's defaults. Let $\mathbf{0}_V$ be the architecture-only model, in which every evidence-set attribute is also reset to
its default. Then, for the pairs that pytm reports on either model:
\[
\mathrm{class}(t,e) =
\begin{cases}
\conf & \text{if } \forall \tau \in \{0,1\}^U: \varphi_t(\sigma,\tau) \\
\ask & \text{if } \varphi_t(\sigma,\mathbf{0}) \wedge \exists \tau: \neg\varphi_t(\sigma,\tau)\\
\ruled & \text{if } \varphi_t(\mathbf{0}_V) \wedge \neg\varphi_t(\sigma,\mathbf{0})
\end{cases}
\]
\conf is thus the Kleene-true case of $\varphi_t$ under partial knowledge~\cite{kleene1952}, restricted to boolean
undecided attributes; a threat that reads a non-boolean undecided attribute (a data classification on a flow, a client's TLS
version) is never \conf. Evidence-backed overrides (e.g.\ S3 Block Public Access) are reported as \emph{mitigated}. The
enumeration is exponential in $|U|$, which is bounded by the attributes a single condition reads; we cap it at $2^{12}$
assignments per pair and conservatively leave a pair with more undecided attributes open (\ask).
\todo{report the observed max and mean $|U|$ over the corpus, and how often the cap is hit.} Every \conf finding carries the evidence of the
attributes that decided it and is linked to the path hop it speaks about.

\paragraph{Cloud pack.} pytm's library targets application threats. A pack of \NumPackThreats cloud threats in pytm's own
format, whose conditions read attributes such as \texttt{imdsV1Allowed} and \texttt{hasInstanceRole}, lets the same
classifier speak about path hops (\NumTOnePackFindings pack findings on the synthetic corpus). Because the pack's attributes come from the same graph as the paths, its
agreement with the paths is not independent evidence; we report pytm's library and the pack separately.

%======================================================================
\section{Questions instead of threat lists}\label{sec:questions}
%======================================================================

An open finding $(t,e)$ needs a set $Q_{t,e}$ of (element, control) keys---the undecided attributes, mapped to the element
that can answer them (a flow's copied attributes are asked of its sink; data attributes of the store holding the data).
Answering all of $Q_{t,e}$ decides the finding either way. \sys asks the key that settles the most findings when added to
the keys already chosen, counting findings on an attack path twice, and repeats while some key settles something:
\[
q_{k+1} = \arg\max_{q} \sum_{(t,e) \in S_k(q)} w(t,e), \qquad A_{k+1} = A_k \cup \{q_{k+1}\},
\]
where $S_k(q) = \{(t,e) : Q_{t,e} \not\subseteq A_k,\ Q_{t,e} \subseteq A_k \cup \{q\}\}$ are the findings that $q$
completes and $w(t,e)$ is 2 on a path and 1 elsewhere.
Unlike maximum coverage, the objective (findings whose \emph{whole} key set is asked) is not submodular, so the classic
$(1-1/e)$ guarantee of greedy selection~\cite{nemhauser1978} does not apply, and a finding that needs two new keys at once is
never credited by a single step. \todo{quantify: share of open findings with $|Q|=1$; compare with an exact ILP on the demo
plan.} Questions are phrased in plain language from a fixed table; data questions (``does this bucket hold personal
data?'') are listed separately because they inform findings that yes/no alone cannot settle.

\paragraph{Answers.} Answers live in \texttt{threat-answers.yaml} next to the Terraform, are reviewed like code, and become
\texttt{APP-} evidence records with trust \emph{attested}. They only set attributes the plan leaves undecided (the plan wins),
malformed or unknown entries are reported, and paths are computed before and independently of them.

%======================================================================
\section{Output for people who do not threat-model}\label{sec:people}
%======================================================================

\paragraph{Path cards.} Each path is rendered deterministically as the card in Figure~\ref{fig:card}: a headline with the
asset, its class and the precondition in words; one sentence per hop with the file and line to change; the smallest fix the
what-if proved; and the first questions on the path's elements. An \emph{everyone} version drops file names, resource
addresses and evidence IDs and adds a short glossary.

\paragraph{Bounded AI.} Six agents (path narrator, plain explainer, remediation, framework mapper, finding explainer, patch
author) receive an evidence bundle, have no tools, and return JSON. A validator accepts an annotation only if its schema is
valid, every cited ID exists and was in the bundle, and its claims do not exceed the evidence (no new resources, no stronger
access, no severity escalation, no instructions from untrusted comments)~\cite{owasp-llm01,ms2026threatai}. Accepted
annotations are displayed next to, never merged into, the deterministic artifacts.

\paragraph{Verified remediation.} The patch author proposes an edit to the Terraform; \sys checks it against an edit policy,
applies it to a copy, re-plans offline, and re-runs paths and classification. A patch is accepted only if the path is gone
(or downgraded) and nothing got worse: no new path, removed resource, added or widened relationship, or new confirmed threat.

%======================================================================
\section{Evaluation}\label{sec:eval}
%======================================================================

\todo{Pin the evaluation commit and the run manifest (tool versions). All numbers below are macros generated from
committed results by \texttt{paper/make\_numbers.py}.}

\paragraph{RQ1: Are the paths right?} \emph{Tier 1:} \NumTOneFixtures synthetic fixtures covering IAM wildcards, trust,
resource policies, explicit deny, conditions, \texttt{for\_each}, \texttt{count}, modules, dynamic blocks, unknown values,
and prompt-injection comments; gold labels were committed before the engine ran on each. Path precision
\NumTOnePathPrecision, recall \NumTOnePathRecall (\NumTOnePathTP TP, \NumTOnePathFP FP, \NumTOnePathFN FN); median runtime
\NumTOneMedianMs\,ms. These fixtures were written by the tool's authors and demonstrate consistency, not accuracy.
\emph{Tier 2:} \NumTTwoCases cases from TerraGoat, IAM Vulnerable and CloudGoat~\cite{terragoat,iamvulnerable,cloudgoat} at frozen
commits, labeled per (entry, target) pair from each scenario's documentation: \NumTTwoTP TP, \NumTTwoFN FN, \NumTTwoFP FP,
\NumTTwoTN TN (precision \NumTTwoPrecision, recall \NumTTwoRecall). \todo{second labeler and agreement ($\kappa$).}

\paragraph{RQ2: Does the output cover known vulnerabilities? (vs.\ TerrARA)} We ran TerrARA's nine benchmark
samples~\cite{terrara-repo,cloudgoat,kaimonkey,awsgoat}. TerrARA publishes per-sample counts but not which vulnerabilities it
counted, so we wrote the list for the samples whose walkthroughs document them (CloudGoat \texttt{glue\_privesc}, AWSGoat
modules 1 and 2: \NumTBVulns vulnerabilities) and committed it, with the matching rules, before the first run. A vulnerability
is covered by a \emph{path} through its resources, a \emph{confirmed} threat of its kind on them, a yes/no question
(\emph{asked}) about a control it names on them, or---TerrARA's own criterion---merely a generic threat of its kind.

The first run found \NumTBFirstRunPaths paths on all nine samples and covered \NumTBFirstRunCovered of \NumTBVulns. The misses
were general AWS semantics missing from the engine: unauthenticated API Gateway routes as entries, ECS services and their
task and host roles, Glue-job PassRole, writing a policy onto a role one can take over, and quoted keys in
\texttt{jsonencode}. After adding them---with Tier~1 and Tier~2 results unchanged---\sys covers \NumTBDeveloper of
\NumTBVulns: \NumTBPath by an attack path matching the walkthrough and \NumTBAsked by a question on the right resource
(\NumTBFirstScreen on the first screen). It asks the reader to look at \NumTBItems items (\NumTBItemsPerVuln per covered
vulnerability; \NumTBItemsFirst on the first screen), against \NumTBTerraraThreats threats from TerrARA
(\NumTBTerraraPerVuln per covered vulnerability) and \NumTBScannerFindings scanner findings. Application bugs (SQL injection,
XSS, IDOR) are invisible in any IaC; the best an IaC tool can do is point at the element and control, which is what the
questions do. \todo{Re-run TerrARA itself (Memgraph + Docker) on the same commits for a like-for-like count.}

\paragraph{RQ3: How much does the plan decide?} On the synthetic corpus, the architecture-only pytm model reports
\NumTOnePytmOnly threats. Fed with plan evidence: \NumTOnePytmDecided decided by the plan, \NumTOnePytmRuledOut ruled out, and
\NumTOnePytmAppTeam open for the app team. On the demo plan, six questions settle \NumDemoSixSettle of \NumDemoOpen open
findings. \todo{distribution over fixtures and Tier 2 of questions needed to settle 50\%/90\%.}

\paragraph{RQ4: Does the AI boundary hold?} Negative tests (fabricated IDs, unsupported claims, escalation, injected
instructions), a prompt-injection fixture, and a pytm model of the AI layer itself. \todo{live-model numbers (two providers,
recorded cassettes), not only the simulated backend: \NumTOneAIGenerated generated, \NumTOneAIAccepted accepted,
\NumTOneAIHeld held.}

\paragraph{RQ5: Are patches proven?} Over \NumTOneRemPaths paths, \NumTOneRemBreak verified breaks and \NumTOneRemRejected
rejected patches; the static what-if agreed with the re-planned result \NumTOneRemWhatIfAgree/\NumTOneRemWhatIfChecked times.
\todo{live patch author.}

\paragraph{RQ6: Can developers act on it?} A within-subjects study (\texttt{paper/study/}) compares three presentations of
the same plan---scanner findings, a pytm threat list, and \sys's card with questions---on whether participants identify the
path, the line to change and the right fix, and how long it takes. \todo{run the study (target $n \ge 24$, half without
security training); report accuracy, time, confidence and SUS~\cite{brooke1996sus}.}

%======================================================================
\section{Discussion and threats to validity}
%======================================================================

\textbf{Construct.} ``Covered by a question'' is weaker than ``covered by a path'': a question directs attention, it does not
show that the bug exists. We report the categories separately and never add them silently. Items to read differ in kind
(a path, a threat, a question); counts are indicative. \textbf{Internal.} The Tier~1 fixtures and the TerrARA ground truth
were written by the authors; both were committed before the runs they score, and the TerrARA engine changes are reported with
the first-run result. \textbf{External.} AWS only; Terraform only; classification needs tags (or declared targets);
application logic is out of reach by design. \textbf{Reliability.} Plans, scanner baselines and results are committed and
the deterministic artifacts are byte-identical across runs.

%======================================================================
\section{Related work}
%======================================================================

\emph{Threat modeling from IaC.} TerrARA~\cite{tran2025terrara} extracts an enriched DFD and elicits threats with
SPARTA~\cite{sion2018sparta}; its limitations section names implicit relationships in policies and threat prioritization as
open problems, which the IAM evaluation and the classification address. StartLeft and IriusRisk import IaC into a rules
engine~\cite{startleft,iriusrisk-terraform}. pytm~\cite{pytm} is threat modeling as code; agent skills now build pytm models
interactively~\cite{tmskills}, complementary to the CI setting here. \emph{Attack graphs.} MulVAL and earlier model-checking
approaches~\cite{ou2005mulval,sheyner2002attackgraphs} compute multi-step attacks over hosts and vulnerabilities; cloud IAM
graph tools (PMapper, Cloudsplaining~\cite{pmapper,cloudsplaining}) analyze deployed accounts. We compute paths before
deployment, from the plan, and attach them to source lines. \emph{IaC analysis.} Security smells~\cite{rahman2019sevensins}
and policy scanners~\cite{checkov,trivy} are per-resource. \emph{Developer-centred security.} Developers act on what their
tools put in front of them~\cite{acar2016sources,tahaei2019survey}; our cards and questions are designed for that.
\emph{LLMs and threat modeling.} Microsoft's guidance on threat modeling AI applications~\cite{ms2026threatai} motivates
keeping the model out of the decision path. \todo{survey 2025--26 LLM-for-threat-modeling papers; verify each citation.}

%======================================================================
\section{Conclusion}
%======================================================================

The plan already knows most of what a threat model needs about the infrastructure. Feeding that evidence into an existing
threat library turns hundreds of generic threats into a few paths, a few confirmed threats, and a few questions---and the
questions are ones a developer can answer. \todo{artifact availability statement (anonymized repository).}

\bibliographystyle{ACM-Reference-Format}
\bibliography{refs}

\end{document}
